% Experimental results for different functionality-driven tasks under three settings: inquiring only documentation, only source code, or both.

\begin{table*}[t]
\centering
\caption{Performance comparison on functionality-driven tasks between the standalone RepoAgent and RepoAgent augmented with source code (RepoAgent + Code).
\textbf{Top XX} indicates the token size of the retrieved documentation and source code context, which is evenly allocated to each source. The largest value in each column is marked in \textbf{bold} and \underline{underlined}.}
\vspace{-1em}

\label{tab:rq2_results}
\setlength{\tabcolsep}{5.5pt}
% \renewcommand{\arraystretch}{0.7}

\begin{tabular}{l|cc|cc|cc|cc|cc|cc}
\toprule
\multirow{2}{*}{\textbf{Model}} & \multicolumn{6}{c|}{\textbf{GPT-4.1}} & \multicolumn{6}{c}{\textbf{Gemini-2.5-pro}} \\
& \multicolumn{2}{c}{\textbf{Top 1024}} & \multicolumn{2}{c}{\textbf{Top 2048}} & \multicolumn{2}{c|}{\textbf{Top 4096}} & \multicolumn{2}{c}{\textbf{Top 1024}} & \multicolumn{2}{c}{\textbf{Top 2048}} & \multicolumn{2}{c}{\textbf{Top 4096}} \\ \midrule
\rowcolor{mygreen} \multicolumn{13}{c}{\textbf{Functionality Detection}} \\ \midrule
\rowcolor{gray!20} \textbf{Metric (\%)} & B-ACC & MCC & B-ACC & MCC & B-ACC & MCC & B-ACC & MCC & B-ACC & MCC & B-ACC & MCC \\ \midrule
RepoAgent  & 62.19 & 25.92 & 62.97 & 26.74 & 63.75 & 27.77 & 63.28 & 25.04 & 67.66 & 33.38 & 69.53 & 37.23 \\
% Code      & \underline{\textbf{68.28}} & \underline{\textbf{36.29}} & \underline{69.06} & \underline{37.04} & \underline{71.09} & \underline{40.23} & \underline{\textbf{73.28}} & \underline{\textbf{44.46}} & \underline{\textbf{74.84}} & \underline{\textbf{47.22}} & \underline{75.16} & \underline{47.86} \\
RepoAgent+Code   & \underline{\textbf{67.81}} & \underline{\textbf{34.80}} & \underline{\textbf{70.47}} & \underline{\textbf{39.33}} & \underline{\textbf{72.66}} & \underline{\textbf{43.27}} & \underline{\textbf{68.13}} & \underline{\textbf{34.32}} & \underline{\textbf{74.53}} & \underline{\textbf{43.52}} & \underline{\textbf{76.25}} & \underline{\textbf{51.89}} \\ \midrule

\rowcolor{myblue} \multicolumn{13}{c}{\textbf{Functionality Localization}} \\ \midrule
\rowcolor{gray!20} \textbf{Metric (\%)} & F1 & IoU & F1 & IoU & F1 & IoU & F1 & IoU & F1 & IoU & F1 & IoU \\ \midrule
RepoAgent & 63.64 & 60.21 & 67.43 & 64.04 & 70.19 & 66.28 & 67.48 & 64.61 & 69.83 & 66.71 & 71.11 & 67.52 \\
% Code & \underline{70.89} & \underline{67.32} & \underline{71.04} & \underline{67.98} & \underline{73.53} & \underline{69.97} & \underline{73.82} & \underline{70.43} & \underline{74.76} & \underline{71.48} & \underline{75.40} & \underline{71.77} \\
RepoAgent+Code & \underline{\textbf{73.09}} & \underline{\textbf{69.62}} & \underline{\textbf{76.65}} & \underline{\textbf{72.94}} & \underline{\textbf{78.22}} & \underline{\textbf{74.58}} & \underline{\textbf{74.49}} & \underline{\textbf{71.03}} & \underline{\textbf{77.75}} & \underline{\textbf{74.39}} & \underline{\textbf{80.39}} & \underline{\textbf{76.98}}\\ \midrule

\rowcolor{mypink} \multicolumn{13}{c}{\textbf{Functionality Completion}} \\ \midrule
\rowcolor{gray!20} \textbf{Metric (\%)} & $\text{EM}_{1.0}$ & $\text{EM}_{0.8}$ & $\text{EM}_{1.0}$ & $\text{EM}_{0.8}$ & $\text{EM}_{1.0}$ & $\text{EM}_{0.8}$ & $\text{EM}_{1.0}$ & $\text{EM}_{0.8}$ & $\text{EM}_{1.0}$ & $\text{EM}_{0.8}$ & $\text{EM}_{1.0}$ & $\text{EM}_{0.8}$ \\
\midrule
RepoAgent & \underline{\textbf{26.40}} & \underline{\textbf{29.37}} & 26.06 & 28.53 & 26.87 & 29.72 & 29.40 & \underline{\textbf{33.20}} & 29.76 & 32.95 & 30.05 & 33.88 \\
% Code     & \underline{\textbf{27.10}} & \underline{\textbf{29.47}} & \underline{26.43} & \underline{29.31} & \underline{27.97} & \underline{30.68} & \underline{\textbf{29.65}} & 32.94 & \underline{30.24} & \underline{34.18} & \underline{30.91} & \underline{34.55} \\
RepoAgent+Code  & 26.22 & 29.11 & \underline{\textbf{26.61}} & \underline{\textbf{30.62}} & \underline{\textbf{28.41}} & \underline{\textbf{31.88}} & \underline{\textbf{29.50}} & 33.13 & \underline{\textbf{30.73}} & \underline{\textbf{34.36}} & \underline{\textbf{31.31}} & \underline{\textbf{35.34}}\\
\bottomrule
\end{tabular}
% \vspace{-1em}
\end{table*}

\subsection{Complementary Value of Source Code}

We further evaluate the complementary value of source code to software documentation on repository comprehension. We design two settings: (1) inquiring only the documentation generated by the best-performance method (RepoAgent),  (2) inquiring both documentation and source code, with code segmented and embedded by syntax structure and context windows evenly allocated. 

As shown in Table~\ref{tab:rq2_results}, the combined approach consistently outperforms inquiring documentation alone. Specifically, it achieves average relative improvements of 10.39\% and 40.35\% for functionality detection, 12.43\% and 12.88\% for functionality localization, and 2.52\% and 3.62\% for functionality completion. These results highlight that source code, by providing direct implementation details, is essential for enhancing repository-level comprehension. Besides, we find that this complementary value is task-dependent. The synergy between documentation and code is most pronounced in functionality detection and localization, with maximum absolute improvements of 8.91\% and 15.50\% for detection, and up to 9.45\% and 9.46\% for localization. This effectiveness stems from the code's ability to supply precise information. However, this synergy is less evident for functionality completion. For instance, under the GPT-4.1 model and Top 1024 tokens context, performance drops slightly by 0.18\% and 0.26\% in $\text{EM}_{1.0}$ and $\text{EM}_{0.8}$, respectively. This indicates that, within a brief overview, broader global information from documentation is more valuable for accurate completion.

% Across two models and three context window settings, this combined approach yields performance gains in 12, 12, and 9 out of the 12 evaluation metrics for three QA tasks. 
% \begin{myfindbox}
%     \textbf{Finding 4:} Source code can offer essential complementary value to software documentation on repository comprehension, especially in functionality detection and localization.
% \end{myfindbox}

\subsection{Implication of Findings}

\subsubsection{Implications for Developers}
Developers should regard documentation as a fundamental knowledge source within their development process and utilize automated tools to enhance documentation generation efficiency. It is advisable to prioritize tools that produce fine-grained and context-rich documentation, and to supplement documentation reading with source code for a deeper repository understanding. An effective strategy is to first consult the documentation for a high-level overview and identification of relevant files, followed by detailed code inspection. This strategy also proves effective during issue resolution, helping developers locate and address issues.

\subsubsection{Implications for Researchers}
Current automated documentation generation methods exhibit notable limitations, particularly in providing intent-oriented details that require a global understanding of functionality. Hence, efforts should be directed toward improving documentation’s ability to provide comprehensive implementation details. Besides, integrating global semantic context proves effective, suggesting that exploring diverse strategies for semantic fusion is a promising direction. Further research should also investigate the broader value of documentation across various development scenarios, such as code review and refactoring, to fully uncover its impact throughout the software development process.


\subsection{Threats and Limitations}
One threat is that \benchmark is limited to 12 popular open-source repositories, which may affect the generalizability of our findings. However, our data construction pipeline is extensible, and we intend to incorporate more repositories in the future. Another threat arises from the inherent randomness of LLMs. Since we use LLMs to answer QA tasks, results may vary across trials. To mitigate this, we conduct multiple runs and report the average results.

% \textbf{Method Selection:} Our experimental analysis is based on six selected methods, which may influence the scope of our findings. To the best of our knowledge, these methods represent the widely adopted and academically recognized approaches available. As new techniques emerge, we will be glad to expand our study to maintain a comprehensive evaluation.

